\subsection{与仿射 Springer 纤维的关系}\label{ss:orb ASF}
从本小节起，限制在 $F$ 为局部函数域的情形，即 $F=k\lr{t}$，其中 $k=\FF_{q}$ 是有限域。设 $\gamma\in\frg(F)$ 正则半单。仿射 Grassmann 簇与仿射 Springer 纤维 $\sX_{\gamma}$ 的定义（\S\ref{s:ASF} 中给出）在基域 $k$ 为有限域时仍有意义，于是得到 $\Gr_{G}$ 的子归纳概形 $\sX_{\gamma}$，两者都定义在 $k$ 上。

下述引理由定义即可看出。

\begin{lemma}\label{l:X gamma k} $k$-有理点集 $\sX_{\gamma}(k)$ 与 \S\ref{sss:X gamma} 中定义的集合 $X_{\gamma}$ 相同，两者都作为 $\Gr_{G}(k)=G(F)/G(\calO_{F})$ 的子集。
\end{lemma}

\subsubsection{} 把基域从 $k$ 变到 $\kbar$，由定理 \ref{th:ft quot} 知 $\L_{\gamma}\bs \sX_{\gamma,\kbar}$ 是 $\kbar$ 上的真代数空间。那里的证明实际上说明此代数空间定义在 $k$ 上，仍记为 $\L_{\gamma}\bs \sX_{\gamma}$。这里要强调：$\L_{\gamma}$ 被看作 $k$ 上的 étale 群概形，其 $\kbar$-点就是此前称为 $\L_{\gamma}$ 的普通群。

由引理 \ref{l:Ogamma} 与引理 \ref{l:X gamma k}，自然会期望轨道积分 $O_{\gamma}(\one_{\frg(\calO_{F})})$ 能表为商 $\L_{\gamma}\bs \sX_{\gamma}$ 上 $k$-点的个数。这一关系在 $G=\GL_{n}$ 时形式最简洁。

\subsubsection{$\GL_{n}$ 的情形}\label{sss:GLn orb coho}
在 \S\ref{sss:GLn orb formula} 的情形中，选取素元 $t_{i}\in F_{i}$ 后定义了格 $L_{0}=t^{\ZZ}_{1}\times \cdots\times t^{\ZZ}_{m}$。把基域从 $k$ 变到 $\kbar$，可以用同样的素元 $t_{i}\in F_{i}$，按 \S\ref{sss:GLn centralizer} 类似地定义格 $\wt\L_{\gamma}$（注意 $F_{i}\wh\otimes_{k}\kbar$ 可能分裂为域之积，但 $t_{i}\otimes 1$ 会投影为每个因子中的素元）。$\Gk$ 在 $\wt\L_{\gamma}$ 上的作用给它 $k$ 上 étale 群概形的结构，正如 $\L_{\gamma}$。有 $\wt\L_{\gamma}(k)=L_{0}$。把定理 \ref{th:ft quot} 中的 $\L_{\gamma}$ 换成 $\wt\L_{\gamma}$，有类似的陈述。特别地，$\wt\L_{\gamma}\bs\sX_{\gamma}$ 是 $k$ 上的真代数空间，且有射影概形到它的满射。

\begin{theorem}\label{th:GLn orb} 设 $G=\GL_{n}$，$\gamma$ 是 $\frg(F)$ 中的正则半单元。把 $G_{\gamma}(F)$ 上的 Haar 测度归一化，使其极大紧子群的体积为 $1$。则
\begin{equation}\label{GLn orb}
O_{\gamma}(\one_{\frg(\calO_{F})})=\#(\wt\L_{\gamma}\bs\sX_{\gamma})(k)=\sum_{i}(-1)^{i}\Tr\left(\Frob_{k}, \cohog{i}{\wt\L_{\gamma}\bs\sX_{\gamma,\kbar},\Ql}\right).
\end{equation}
\end{theorem}

\eqref{GLn orb} 中的第二个等号来自 Grothendieck--Lefschetz 迹公式。与 \eqref{GLn L0} 比较，只需说明 $L_{0}\bs X_{\gamma}=\wt\L_{\gamma}(k)\bs \sX_{\gamma}(k)$ 与 $(\wt\L_{\gamma}\bs\sX_{\gamma})(k)$ 相同。先对商叠的 $k$-点作一些一般性说明。

\subsubsection{商的 $k$-点}\label{sss:k pts quot} 考虑商叠 $\sY=[A\bs X]$，其中 $X$ 是 $k$ 上概形，$A$ 是 $k$ 上作用在 $X$ 上的代数群。按定义，$\sY(k)$ 是由数对 $(S, f)$ 构成的群胚，其中 $S\to \Spec k$ 是左 $A$-主齐性空间，$f:S\to X$ 是 $A$-等变态射。$A$-主齐性空间 $S$ 的同构类由 Galois 上同调 $\cohog{1}{k, A}:=\cohog{1}{\Gk,A(\kbar)}$ 分类。对每个类 $\xi\in \cohog{1}{k,A}$，设 $S_{\xi}$ 是类为 $\xi$ 的 $k$ 上 $A$-主齐性空间。可以用 $X_{\xi}:=A\bs (S_{\xi}\times X)$ 定义 $X$ 在 $k$ 上的扭形式。还有作用在 $k$-概形 $X_{\xi}$ 上的内形式 $A_{\xi}:=\Aut_{A}(S_{\xi})$。容易看出 $A$-等变态射 $f:S_{\xi}\to X$ 与 $X_{\xi}(k)$ 一一对应。因此得到群胚分解
\begin{equation}\label{k pts quot}
\sY(k)=[A\bs X](k)\cong\bigsqcup_{\xi\in\cohog{1}{k,A}}A_{\xi}(k)\bs X_{\xi}(k).
\end{equation}

\subsubsection{定理 \ref{th:GLn orb} 的证明} 已化归为证明 $(\wt\L_{\gamma}\bs\sX_{\gamma})(k)= \wt\L_{\gamma}(k)\bs \sX_{\gamma}(k)$。由 \eqref{k pts quot}，只需说明 $\cohog{1}{k,\wt\L_{\gamma}}=0$。沿用 \S\ref{sss:GLn orb coho} 的记号。由 $\wt\L_{\gamma}$ 的定义可见，作为 $\Gk$-模它有形式
\begin{equation*}
\wt\L_{\gamma}\cong \bigoplus_{i=1}^{m}\Ind^{\Gk}_{\Gal(\kbar/k_{i})}\ZZ,
\end{equation*}
其中 $k_{i}$ 是 $F_{i}$ 的剩余域。因此
\begin{equation*}
\cohog{1}{k,\wt\L_{\gamma}}\cong\bigoplus_{i=1}^{m}\cohog{1}{k_{i},\ZZ}=0.
\end{equation*}
\qed

\subsection{稳定轨道积分} 设定与 \S\ref{ss:orb ASF} 相同。对一般的 $G$，公式 \eqref{GLn orb} 的推广并不直接。具体说来，轨道积分 $O_{\gamma}(\one_{\frg(\calO_{F})})$ 本身并没有上同调解释。问题在于未必能找到 $\wt\L_{\gamma}$ 的类比物、使其一阶 Galois 上同调消没，从而 $(\wt\L_{\gamma}\bs\sX_{\gamma})(k)$ 就简单地等于 $\wt\L_{\gamma}(k)\bs \sX_{\gamma}(k)$。由公式 \eqref{k pts quot}，一个自然的补救办法是把 $\sX_{\gamma}$ 的各扭形式放在一起考虑。这提示我们不要只看轨道积分 $O_{\gamma}(\one_{\frg(\calO_{F})})$，而要看若干个轨道积分 $O_{\gamma'}(\one_{\frg(\calO_{F})})$ 之和。

\subsubsection{稳定共轭} 固定正则半单元 $\gamma\in\frg(F)$。称 $\gamma'\in\frg(F)$ {\em 稳定共轭}于 $\gamma$，若它与 $\gamma$ 在同一 $G(\ov{F})$-轨道中。等价地，$\gamma'$ 稳定共轭于 $\gamma$，当且仅当 $\chi(\gamma')=\chi(\gamma)\in\frc(F)$。对稳定共轭于 $\gamma$ 的 $\gamma'$，可以附上一个 Galois 上同调类 $\inv(\gamma,\gamma')\in\cohog{1}{F,G_{\gamma}}$，它在 $\cohog{1}{F,G}$ 中的像平凡。对应 $\gamma'\mapsto\inv(\gamma,\gamma')$ 给出带点集的双射
\begin{equation}\label{stable class}
\{\gamma'\in\frg(F) \mbox{ 稳定共轭于 } \gamma\}/G(F) \cong \ker(\cohog{1}{F,G_{\gamma}}\to \cohog{1}{F,G}).
\end{equation}

\subsubsection{$\GL_{n}$ 的情形} 当 $G=\GL_{n}$ 时，沿用 \S\ref{sss:GLn centralizer} 的记号。有
\begin{equation*}
\cohog{1}{F,G_{\gamma}}\cong \prod_{i=1}^{m}\cohog{1}{F_{i},\Gm}=\jiao{1}.
\end{equation*}
因此由 \eqref{stable class}，所有稳定共轭于 $\gamma$ 的元素实际上都在 $G(F)$ 中共轭于 $\gamma$。当然，此陈述也可以用友矩阵直接证明。

\subsubsection{$\SL_{2}$ 的情形}\label{sss:st conj SL2} 考虑 $G=\SL_{2}$。设 $a\in k^{\times}-(k^{\times})^{2}$。令
\begin{equation*}
\gamma=\mat{0}{at}{t}{0}, \quad \gamma'=\mat{0}{at^{2}}{1}{0}
\end{equation*}
为 $\sl_{2}(F)$ 中两个正则半单元。由于它们的行列式都是 $-at^{2}$，两者稳定共轭。然而它们在 $\SL_{2}(F)$ 下并不共轭。可以证明 $\gamma$ 的稳定共轭类恰由两个 $\SL_{2}(F)$-轨道组成，分别以 $\gamma$ 与 $\gamma'$ 为代表，见习题 \ref{exIII:stable conj}。

\subsubsection{稳定轨道积分的定义} 设 $\varphi\in\sS(\frg(F))$。定义 $\varphi$ 关于 $\gamma$ 的{\em 稳定轨道积分}为
\begin{equation}\label{SO}
SO_{\gamma}(\varphi)=\sum_{\gamma'}O_{\gamma'}(\varphi),
\end{equation}
其中 $\gamma'$ 遍历 $\frg(F)$ 中稳定共轭于 $\gamma$ 的元素的 $G(F)$-轨道。对稳定共轭于 $\gamma$ 的 $\gamma'$，有 $F$-群之间的{\em 典范同构} $G_{\gamma}\cong G_{\gamma'}$。因此一旦固定 $G_{\gamma}(F)$ 上的 Haar 测度，就得到其他中心化子 $G_{\gamma'}(F)$ 上的典范 Haar 测度。\eqref{SO} 中的 $O_{\gamma'}(\varphi)$ 正是按此选取来定义的。

\subsubsection{上同调的稳定部分}\label{sss:st H} 商群概形 $LG_{\gamma}/\L_{\gamma}$ 作用在 $\L_{\gamma}\bs\sX_{\gamma}$ 上。分支群 $\pi_{0}(LG_{\gamma}/\L_{\gamma})$ 是 $k$ 上 étale 群概形，其 $\kbar$-点作用在 $\cohog{*}{L_{\gamma}\bs\sX_{\gamma,\kbar},\Ql}$ 上。定义此上同调群的{\em 稳定部分}为此作用下的不变量
\begin{equation*}
\cohog{*}{\L_{\gamma}\bs\sX_{\gamma,\kbar}}_{\st}=\cohog{*}{\L_{\gamma}\bs\sX_{\gamma,\kbar}}^{\pi_{0}(LG_{\gamma}/\L_{\gamma})(\kbar)}.
\end{equation*}
可以证明：若把 $\L_{\gamma}$ 换成 $G_{\gamma}(F^{\ur})$ 中任一与 $\L_{\gamma}$ 可公度的 $\Gk$-稳定自由交换子群 $\L$，类似定义的稳定部分上同调与上述典范同构。

\begin{theorem}[Goresky--Kottwitz--MacPherson {\cite[定理 15.8]{GKM}} 与 Ng\^{o} {\cite[推论 8.2.10]{NgoFL}} 的特例]\label{th:orb coho} 有
\begin{equation}\label{d}
SO_{\gamma}(\one_{\frg(\calO_{F})})=\frac{1}{\vol(K_{\gamma},\mu_{G_{\gamma}})}\sum_{i}(-1)^{i}\Tr\left(\Frob_{k}, \cohog{i}{\L_{\gamma}\bs\sX_{\gamma,\kbar},\Ql}_{\st}\right).
\end{equation}
这里 $K_{\gamma}\subset G_{\gamma}(F)$ 是环面 $G_{\gamma}$ 的 parahoric 子群。
\end{theorem}

稳定轨道积分的这一上同调解释，是基本引理之证明的出发点（见 \cite{GKM} 与 \cite{NgoFL}）。

\subsubsection{}\label{sss:K gamma} 简略说明 parahoric 子群 $K_{\gamma}$ 的定义。令 $\tilK_{\gamma}\subset G_{\gamma}(F)$ 为极大紧子群。则存在 $\calO_{F}$ 上典范光滑群概形 $\calG_{\gamma}$，其 $F$-纤维为 $G_{\gamma}$，其 $\calO_{F}$-点为 $\tilK_{\gamma}$。此群概形 $\calG_{\gamma}$ 就是环面 $G_{\gamma}$ 的有限型 Néron 模型。令 $\calG^{\circ}_{\gamma}\subset\calG_{\gamma}$ 为去掉 $\calG_{\gamma}$ 特殊纤维的非中性分支后得到的开子群概形。则按定义 $K_{\gamma}=\calG^{\circ}_{\gamma}(\calO_{F})$。正环路群 $L^{+}\calG^{\circ}_{\gamma}$ 在相差幂零元的意义下是 $LG_{\gamma}$ 的中性分支。

\subsubsection{定理 \ref{th:orb coho} 的证明梗概}\label{sss:pf SO} 下面概述的论证更接近 \cite{GKM} 而非 \cite{NgoFL}。可以证明存在 étale $k$-子群 $\wt\L\subset LG_{\gamma}$，它与 $\L_{\gamma}$ 可公度，且{\em 满射}到 $k$ 上 étale 群概形 $\pi_{0}(LG_{\gamma})$（其核必有限）。群 $\wt\L$ 是 \S\ref{sss:GLn orb coho} 中定义的 $\wt\L_{\gamma}$ 的类比物。作商叠 $[\wt\L\bs \sX_{\gamma}]$，则有同构
\begin{equation}\label{st coho quot}
\cohog{*}{\L_{\gamma}\bs\sX_{\gamma,\kbar},\Ql}_{\st}\cong \cohog{*}{[\wt\L\bs\sX_{\gamma,\kbar}],\Ql}.
\end{equation}
事实上，由 \S\ref{sss:st H} 末尾的讨论，把 $\L_{\gamma}$ 换成可公度的格时上式左端不变，故可把 $\L_{\gamma}$ 缩小，假定 $\L_{\gamma}\subset \wt\L$。另一方面，上式右端可用与映射 $\L_{\gamma}\bs\sX_{\gamma,\kbar}\to[\wt\L\bs\sX_{\gamma,\kbar}]$ 相联系的 Leray 谱序列计算；此映射是有限离散群 $(\wt\L/\L_{\gamma})(\kbar)$ 下的主齐性空间，因此 $\cohog{*}{[\wt\L\bs\sX_{\gamma,\kbar}],\Ql}=\cohog{*}{\L_{\gamma}\bs\sX_{\gamma,\kbar},\Ql}^{(\wt\L/\L_{\gamma})(\kbar)}$。由于按 $\wt\L$ 的选取，$(\wt\L/\L_{\gamma})(\kbar)$ 满射到 $\pi_{0}(LG_{\gamma}/\L_{\gamma})(\kbar)$，便有
\begin{equation*}
\cohog{*}{\L_{\gamma}\bs\sX_{\gamma,\kbar},\Ql}^{\pi_{0}(LG_{\gamma}/\L_{\gamma})(\kbar)}=\cohog{*}{\L_{\gamma}\bs\sX_{\gamma,\kbar},\Ql}^{(\wt\L/\L_{\gamma})(\kbar)},
\end{equation*}
由此即得 \eqref{st coho quot}。

由 \S\ref{sss:k pts quot} 与 \eqref{k pts quot} 的讨论，有
\begin{equation}\label{L bs X decomp}
[\wt\L\bs\sX_{\gamma,\kbar}](k)\cong\bigsqcup_{\xi\in\cohog{1}{k,\wt\L}} \wt\L(k)\bs\sX_{\gamma,\xi}(k).
\end{equation}
包含 $\wt\L\subset LG_{\gamma}$ 给出
\begin{equation*}
\theta: \cohog{1}{k,\wt\L}\surj\cohog{1}{k,LG_{\gamma}}\cong\cohog{1}{F, G_{\gamma}}.
\end{equation*}
第一个满射来自 Lang 定理，因为商 $LG_{\gamma}/\wt\L$ 连通；第二个来自 Lang 的另一定理，它断言 $\cohog{1}{F^{\ur}, G_{\gamma}}$ 消没\footnote{参见 \cite[Ch.X, \S7, p.170, Application and Example (b)]{Serre}。设 $K$ 是剩余域完全的完备离散赋值域，$K^{\ur}$ 为其极大非分歧扩张。Lang 定理断言 $K^{\ur}$ 是 $C_{1}$-域。因此对 $K$ 上任一环面 $A$ 有 $\cohog{1}{K^{\ur}, A}=0$。}。对每个满足 $\sX_{\gamma,\xi}(k)\neq\varnothing$ 的 $\xi\in\cohog{1}{k,\wt\L}$，可以证明 $\theta(\xi)$ 在 $\cohog{1}{F,G}$ 中的像平凡。因此由 \eqref{stable class}，可以对每个 $\xi\in\cohog{1}{k,\wt\L}$ 附上一个稳定共轭于 $\gamma$ 的元素 $\gamma'$，它在 $G(F)$-共轭下唯一，且满足 $\inv(\gamma,\gamma')=\theta(\xi)$。可以证明 $\sX_{\gamma,\xi}(k)$ 与集合 $X_{\gamma'}$ 一一对应。因此 \eqref{L bs X decomp} 蕴含
\begin{equation}\label{pre SO}
\#[\wt\L\bs\sX_{\gamma}](k)=\#\ker(\theta)\sum_{\gamma'}\#(\wt\L(k)\bs X_{\gamma'}),
\end{equation}
其中求和遍历稳定共轭于 $\gamma$ 的那些 $\gamma'$ 的 $G(F)$-轨道。把引理 \ref{l:Ogamma} 用于离散余紧子群 $\wt\L(k)$，得
\begin{equation*}
O_{\gamma'}(\one_{\frg(\calO_{F})})=\frac{1}{\vol(G_{\gamma}(F)/\wt\L(k),\mu_{G_{\gamma}})}\#(\wt\L(k)\bs X_{\gamma'}).
\end{equation*}
代入 \eqref{pre SO} 的右端，得到
\begin{equation}\label{next SO}
\SO_{\gamma}(\one_{\frg(\calO_{F})})=\frac{1}{\#\ker(\theta)\vol(G_{\gamma}(F)/\wt\L(k),\mu_{G_{\gamma}})}\#[\wt\L\bs\sX_{\gamma}](k).
\end{equation}
由 Grothendieck--Lefschetz 迹公式，$\#[\wt\L\bs\sX_{\gamma,\kbar}](k)$ 等于 $\cohog{*}{[\wt\L\bs\sX_{\gamma,\kbar}],\Ql}$ 上的交错 Frobenius 迹；由 \eqref{st coho quot}，后者可以等同于 $\cohog{*}{\L_{\gamma}\bs\sX_{\gamma,\kbar},\Ql}_{\st}$。因此定理归结为等式 \eqref{next SO} 与如下体积等式
\begin{equation}\label{vol id}
\#\ker(\theta)\vol(G_{\gamma}(F)/\wt\L(k),\mu_{G_{\gamma}})=\vol(K_{\gamma},\mu_{G_{\gamma}}).
\end{equation}
为证此式，令 $C=L^{+}\calG^{\circ}_{\gamma}\cap \wt\L$（其中 $\calG^{\circ}_{\gamma}$ 是 $G_{\gamma}$ 的连通 Néron 模型，其 $\calO_{F}$-点为 $K_{\gamma}$，见 \S\ref{sss:K gamma}）。这是 $k$ 上有限 étale 群。有 $k$ 上群归纳概形的短正合列
\begin{equation*}
1\to C\to L^{+}\calG^{\circ}_{\gamma}\times\wt\L\to (LG_{\gamma})^{\red}\to 1.
\end{equation*}
与之相伴的 $\Gk$-上同调六项正合列给出
\begin{equation*}
1\to C(k)\to K_{\gamma}\times \wt\L(k)\to G_{\gamma}(F)\to \cohog{1}{k,C}\to \ker(\theta)\to1,
\end{equation*}
由此并利用 $\#C(k)=\#\cohog{1}{k,C}$，即得 \eqref{vol id}。\qed
